\documentclass[11pt]{article}
\usepackage{amsmath,amssymb,amsthm}
\usepackage{hyperref}

\newtheorem{theorem}{Theorem}
\newtheorem{conjecture}[theorem]{Conjecture}
\newtheorem{remark}[theorem]{Remark}
\newtheorem{observation}[theorem]{Observation}
\newtheorem{question}[theorem]{Question}

\newcommand{\CPP}{\operatorname{CPP}}
\newcommand{\GK}{\operatorname{GK}}
\newcommand{\CT}{\operatorname{CT}}
\newcommand{\QSym}{\operatorname{QSym}}
\newcommand{\Sym}{\operatorname{Sym}}
\newcommand{\qpoch}[2]{(#1;#2)_\infty}
\newcommand{\qpochn}[3]{(#1;#2)_{#3}}
\newcommand{\sln}{\widehat{\mathfrak{sl}}_n}

\title{Literature Review: Constant Term Identities\\
and the $Q$-Polynomial Extraction}
\author{Robin Langer}
\date{February 2026}

\begin{document}
\maketitle

\begin{abstract}
We survey the theory of constant term identities from Dyson
through Macdonald and Cherednik, and propose reframing the
extraction of Warnaar's $Q_{n,c}(q)$ polynomials as a constant
term identity. The $Q$-polynomial is
$Q_{n,c}(q) = (q;q)_n \cdot \CT_z[z^{-n}(zq;q)_\infty
\cdot \GK_c(z,q)]$,
which has the structural form of a Macdonald-type constant term
with the Borodin product playing the role of the root-system
kernel. The Karolyi--Lascoux--Warnaar (2015) extension of
Macdonald's constant term identities with extra parameters is
identified as the most promising point of contact.
\end{abstract}

%% ========================================================================
\section{The constant term lineage}
\label{sec:lineage}
%% ========================================================================

\subsection{Dyson's conjecture and the $q$-Dyson theorem}

Dyson~\cite{dyson-1962} conjectured that the constant term of
the Laurent polynomial
\[
  \prod_{1 \le i \ne j \le n}
  \Bigl(1 - \frac{x_i}{x_j}\Bigr)^{a_i}
\]
equals the multinomial coefficient
$(a_1 + \cdots + a_n)! / (a_1! \cdots a_n!)$.
This was proved by Gunson and Wilson (1962), with a celebrated
combinatorial proof by Zeilberger~\cite{zeilberger-1982}.

Andrews formulated the $q$-analog: the constant term of
\[
  \prod_{1 \le i < j \le n}
  (x_i/x_j)_{a_i}\, (q\, x_j/x_i)_{a_j}
\]
equals $(q)_{a_1 + \cdots + a_n} / \bigl((q)_{a_1} \cdots
(q)_{a_n}\bigr)$. This was proved by Zeilberger and
Bressoud~\cite{zeilberger-bressoud-1985}.

\subsection{The Morris identity and the Selberg integral}

The Selberg integral~\cite{selberg-1944} evaluates a
multivariate beta integral with Vandermonde-squared weight.
It is equivalent to the Morris constant term
identity~\cite{morris-1982}, which extracts the constant term
of a product of Laurent polynomials on the type-$A$ root system.

The $q$-analog (the Habsieger--Kadell $q$-Morris identity) was
proved independently by Habsieger~\cite{habsieger-1988} and
Kadell~\cite{kadell-1988}. The chain of generalizations is:
\[
  \text{Selberg integral}
  \;\iff\; \text{Morris CT}
  \;\xrightarrow{q}
  \text{Habsieger--Kadell}
  \;\hookrightarrow\; \text{Macdonald CT (type $A_n$)}.
\]

\subsection{Macdonald's constant term conjecture}

Macdonald~\cite{macdonald-1982} vastly generalised Dyson's
conjecture to arbitrary reduced root systems~$R$. For each~$R$
with Weyl group~$W$, the conjecture gives a closed-form
evaluation of the constant term of a product over positive
roots, involving parameters tied to the degrees of the
fundamental invariants of~$W$.

Case-by-case proofs were given for the classical types:
\begin{itemize}
  \item Type $A_n$: Zeilberger--Bressoud (the $q$-Dyson theorem).
  \item Types $B_n$, $D_n$: Kadell.
  \item Types $BC_n$, $C_n$: Gustafson.
  \item Exceptional types $G_2$, $F_4$: Zeilberger, Habsieger,
    Garvan--Gonnet.
\end{itemize}

The complete, classification-free proof for all root systems was
achieved by Cherednik~\cite{cherednik-1995} using the theory of
\textbf{double affine Hecke algebras} (DAHAs). Cherednik's proof
bypassed case-by-case verification and provided a unified
algebraic framework.

\subsection{Karolyi--Lascoux--Warnaar (2015): extra parameters}

Karolyi, Lascoux, and Warnaar~\cite{klw-2015} proved that
Macdonald's constant term identities admit an
\textbf{extra set of free parameters}, linking them to
Poincar\'e polynomials of classical root systems. Their proof
combined the Karasev--Petrov polynomial lemma (from the
Combinatorial Nullstellensatz tradition), scalar products for
Demazure characters, and properties of $(0,1)$-matrices.

As an application, they also proved Sills' constant term
conjecture and Kadell's orthogonality conjecture.

\begin{remark}
Note that \textbf{Warnaar} is a co-author on this paper.
The same Warnaar whose Conjecture~2.7 we are studying.
His expertise in extending Macdonald-type constant term
identities with extra parameters makes this connection
especially promising.
\end{remark}

%% ========================================================================
\section{Constant term methods for positivity}
\label{sec:positivity}
%% ========================================================================

The standard pattern for proving positivity via constant terms:

\begin{enumerate}
  \item Express the quantity as
    $\CT[\text{kernel} \cdot \text{generating function}]$.
  \item Show the product evaluates to a closed form
    (via Macdonald/Cherednik/Selberg machinery).
  \item Closed forms are products of $q$-Pochhammer symbols,
    which are manifestly positive.
\end{enumerate}

This pattern succeeds in the following cases.

\subsection{Zeilberger's KOH theorem}

Zeilberger reformulated O'Hara's combinatorial proof of the
unimodality of $q$-binomial coefficients as a constant term
decomposition~\cite{zeilberger-koh}. The $q$-binomial
$\binom{m+n}{n}_q$ is decomposed into a finite sum of
symmetric, unimodal polynomials with non-negative coefficients.
Pak and Panova later proved \emph{strict} unimodality using
Kronecker coefficients.

\subsection{Kronecker coefficients via constant terms}

Rosas and Zabrocki~\cite{rosas-zabrocki} developed a framework
for computing Kronecker coefficients via constant term extraction
from symmetric functions. Since Kronecker coefficients are
multiplicities (hence non-negative), this provides a bridge
between constant term positivity and representation theory.

\subsection{Warnaar's $q$-rious positivity conjecture}

Warnaar and Zudilin~\cite{warnaar-zudilin-2011} posed a
striking positivity conjecture for certain $q$-series involving
$q$-binomials and factorial ratios. The conjecture asserts that
a specific rational function of~$q$ is always a polynomial with
non-negative integer coefficients. This remains \textbf{wide
open} and has the same flavour as Conjecture~2.7: a $q$-series
that ``should'' be positive but resists proof.

%% ========================================================================
\section{Cylindric partitions and constant terms}
\label{sec:cpp}
%% ========================================================================

\subsection{Gessel--Krattenthaler: affine root system formulas}

Gessel and Krattenthaler~\cite{gessel-krattenthaler-1997} showed
that cylindric partition generating functions produce summation
and transformation formulas for basic hypergeometric series
associated to the \textbf{affine root system}
$\widetilde{A}_r$---the same root system family as the
Macdonald constant term identities of type~$A$.

\subsection{Borodin's product formula}

Borodin~\cite{borodin-2007} proved that the generating function
for cylindric partitions of any profile~$c$ is an infinite
product of $q$-Pochhammer symbols---a ``hook-product formula''
on the cylinder. The key structural consequence:
$\GK_c(z,q)$ is always ``nice'' (an infinite product).

\subsection{Postnikov: cylindric Schur positivity}

Postnikov~\cite{postnikov-2005} introduced toric (cylindric)
Schur polynomials and proved that their coefficients in the
Schur expansion are 3-point Gromov--Witten invariants of the
Grassmannian. In particular, all toric Schur polynomials are
\textbf{Schur-positive}. McNamara later extended this to
cylindric skew Schur functions~\cite{mcnamara-2019}.

%% ========================================================================
\section{The $Q$-polynomial as a constant term}
\label{sec:Q-CT}
%% ========================================================================

\subsection{The reframing}

Recall that Warnaar's $Q$-polynomial is
\begin{equation}\label{eq:Q-def}
  Q_{n,c}(q)
  = (q;q)_n \cdot [z^n]\bigl(\qpoch{zq}{q} \cdot \GK_c(z,q)\bigr).
\end{equation}
Rewriting the coefficient extraction as a constant term:
\begin{equation}\label{eq:Q-CT}
  \boxed{
  Q_{n,c}(q)
  = (q;q)_n \cdot
    \CT_z\Bigl[z^{-n} \cdot \qpoch{zq}{q} \cdot \GK_c(z,q)\Bigr].
  }
\end{equation}

\subsection{The analogy}

The classical extraction from an exponential generating function
is $a_n = n! \cdot [x^n] f(x) = \CT_x[x^{-n} \cdot n! \cdot
f(x)]$, where the $n!$ and the $e^{-x}$ factor combine into a
shift operator. The $q$-analogue is:
\begin{center}
\begin{tabular}{c|c}
  \textbf{Classical} & \textbf{$q$-analogue} \\
  \hline
  $n!$ & $(q;q)_n$ \\
  $e^{-x}$ (= $\sum (-x)^k/k!$) & $\qpoch{zq}{q}$
    (= $\prod_{k \geq 1}(1 - zq^k)$) \\
  $[x^n]$ & $\CT_z[z^{-n} \cdot]$ \\
  $f(x)$ & $\GK_c(z,q)$
\end{tabular}
\end{center}

In the symmetric function world:
\begin{itemize}
  \item $\qpoch{zq}{q}$ under the principal specialization
    $x_i = q^i$ is the generating function
    $E(-z) = \sum_k (-z)^k e_k$ for elementary symmetric
    functions.
  \item The fundamental relation $E(-z) \cdot H(z) = 1$
    (where $H(z) = \sum h_n z^n$) ties the extraction to the
    $\omega$-involution.
  \item The alternating structure
    $\sum (-1)^k e_k \cdot h_{n-k}^{\perp}$ is precisely a
    \textbf{Bernstein vertex operator} $B_n$, which creates
    rows of Schur functions: $B_n \cdot s_\lambda = s_{(n,\lambda)}$.
\end{itemize}

\subsection{Structure of the constant term}

Substituting Borodin's product formula, the constant term
becomes:
\begin{equation}\label{eq:Q-CT-explicit}
  Q_{n,c}(q)
  = (q;q)_n \cdot
    \CT_z\Bigl[
      z^{-n} \cdot
      \underbrace{\qpoch{zq}{q}}_{\text{elementary factor}}
      \cdot
      \underbrace{\text{Borodin product}}_{\sln[r] \text{ character}}
    \Bigr].
\end{equation}
The integrand is:
\begin{itemize}
  \item $z^{-n}$: the extraction variable.
  \item $\qpoch{zq}{q} = \prod_{k \geq 1}(1 - zq^k)$:
    a type-$A$ $q$-Pochhammer in the extraction variable.
  \item The Borodin product: a ratio of $q$-Pochhammer symbols
    and theta functions with $\widetilde{A}_r$ structure.
\end{itemize}
This has the structural form of a \textbf{Macdonald-type
constant term}: a product over a root system, with
$z^{-n}$ playing the role of an extra parameter.

%% ========================================================================
\section{Why this might work}
\label{sec:why}
%% ========================================================================

\subsection{The Gessel--Krattenthaler connection}

Cylindric partition generating functions produce formulas
for the affine root system $\widetilde{A}_r$, which is
the natural home of the Macdonald constant term identities
(type~$A$). The $Q$-extraction adds a parameter ($n$, the
max entry) to this framework. If this parameter corresponds
to one of the ``extra free parameters'' in the
Karolyi--Lascoux--Warnaar extension, the constant term would
evaluate to a product and positivity would follow.

\subsection{The $q$-shift structure}

The Corteel--Welsh functional equations are $q$-shift relations:
\begin{equation}\label{eq:CW-shift}
  G_c(y) = \sum_{\emptyset \subset J \subseteq I_c}
  (-1)^{|J|-1}\,
  (yq;q)_{|J|-1}\, G_{c(J)}(yq^{|J|}).
\end{equation}
In the constant term picture, the $q$-shift $y \mapsto yq^{|J|}$
becomes multiplication by $z^{|J|}$ inside the $\CT$.
The functional equations translate into
\textbf{recurrences for the constant term} relating different
profiles and different values of~$n$. If these recurrences
preserve positivity, one could bootstrap from known base cases.

\subsection{Multivariate promotion}

The current formulation uses a single extraction variable~$z$.
But the Borodin product has $\widetilde{A}_r$ structure
(rank~3 in our case), with three rows of partitions
$(\nu^{(1)}, \nu^{(2)}, \nu^{(3)})$. Promoting
$z$ to $(z_1, z_2, z_3)$---one per row---and taking a
\emph{multivariate} constant term might land squarely in the
$q$-Morris / Macdonald world:
\begin{equation}\label{eq:Q-multivariate}
  Q_{n,c}(q)
  \;\stackrel{?}{=}\;
  (q;q)_n \cdot
  \CT_{z_1, z_2, z_3}\Bigl[
    \text{kernel}(z_1, z_2, z_3; q)
    \cdot \text{Borodin}(z_1, z_2, z_3; q)
  \Bigr].
\end{equation}
This multivariate constant term would be a natural candidate
for evaluation via the Cherednik/DAHA machinery.

\begin{question}
Does the multivariate promotion~\eqref{eq:Q-multivariate}
produce a constant term identity that falls into the
Macdonald framework?
\end{question}

%% ========================================================================
\section{Key references to follow up}
\label{sec:refs}
%% ========================================================================

In order of priority:

\begin{enumerate}
  \item \textbf{Karolyi--Lascoux--Warnaar (2015)}~\cite{klw-2015}:
    the extra-parameter extension of Macdonald's constant term
    identities. This is the most likely point of contact.
  \item \textbf{Warnaar--Zudilin (2011)}~\cite{warnaar-zudilin-2011}:
    the ``$q$-rious positivity'' conjecture. Same flavour as
    Conjecture~2.7; understanding why one is hard may illuminate
    the other.
  \item \textbf{Warnaar, ``$q$-Selberg integrals and Macdonald
    polynomials''}~\cite{warnaar-qselberg}: connects $q$-Selberg
    integrals to Macdonald polynomial theory. The $q$-Selberg
    integral is the continuous analogue of the $q$-Morris
    constant term.
  \item \textbf{Gessel--Krattenthaler (1997)}~\cite{gessel-krattenthaler-1997}:
    the original paper establishing the $\widetilde{A}_r$
    structure of cylindric partition generating functions.
  \item \textbf{Cherednik (1995)}~\cite{cherednik-1995}:
    the DAHA proof of Macdonald's conjecture. The machinery
    that would be needed if the multivariate promotion works.
\end{enumerate}

%% ========================================================================
\section{Summary}
\label{sec:summary}
%% ========================================================================

\begin{center}
\begin{tabular}{l|l}
  \textbf{Observation} & \textbf{Implication} \\
  \hline
  $Q_{n,c}$ is a constant term~\eqref{eq:Q-CT}
    & Connects to Dyson--Macdonald--Cherednik lineage \\
  Kernel is $z^{-n}(zq;q)_\infty$
    & Type-$A$ $q$-Pochhammer $\times$ monomial \\
  $\GK_c$ is a product (Borodin)
    & Integrand is a product over $\widetilde{A}_r$ \\
  Gessel--Krattenthaler: CPPs $\leftrightarrow$ $\widetilde{A}_r$
    & Same root system as Macdonald CT \\
  KLW (2015): extra parameters in Macdonald CT
    & The max-entry $n$ could be such a parameter \\
  Corteel--Welsh = $q$-shift relations
    & Translate to CT recurrences (may preserve positivity) \\
  Multivariate promotion $z \to (z_1, z_2, z_3)$
    & Could land in $q$-Morris / DAHA territory
\end{tabular}
\end{center}

\medskip

The reframing of $Q_{n,c}$ as a constant term moves the
positivity question from ``why are these coefficients positive''
to ``why is this constant term positive''---a much
better-studied question with powerful machinery
(Cherednik algebras, $q$-Selberg integrals, the KLW
extension) available. The most promising next step is to check
whether the constant term~\eqref{eq:Q-CT-explicit}, with
Borodin's product substituted explicitly, matches any known
family of Macdonald-type evaluations.

%% ========================================================================
%% References
%% ========================================================================

\begin{thebibliography}{99}

\bibitem{borodin-2007}
A.~Borodin,
Periodic Schur process and cylindric partitions,
\emph{Duke Math.\ J.}\ \textbf{140}(3) (2007), 391--468.
arXiv:math/0601019.

\bibitem{cherednik-1995}
I.~Cherednik,
Double affine Hecke algebras and Macdonald's conjectures,
\emph{Ann.\ of Math.}\ (2) \textbf{141}(1) (1995), 191--216.

\bibitem{corteel-welsh-2019}
S.~Corteel and T.~Welsh,
The $A_2$ Rogers--Ramanujan identities revisited,
arXiv:1905.08343, 2019.

\bibitem{dyson-1962}
F.~J.~Dyson,
Statistical theory of the energy levels of complex systems.~I,
\emph{J.\ Math.\ Phys.}\ \textbf{3} (1962), 140--156.

\bibitem{gessel-krattenthaler-1997}
I.~M.~Gessel and C.~Krattenthaler,
Cylindric partitions,
\emph{Trans.\ Amer.\ Math.\ Soc.}\ \textbf{349}(2) (1997),
429--479.

\bibitem{habsieger-1988}
L.~Habsieger,
Une $q$-int\'egrale de Selberg et Askey,
\emph{SIAM J.\ Math.\ Anal.}\ \textbf{19}(6) (1988),
1475--1489.

\bibitem{kadell-1988}
K.~W.~J.~Kadell,
A proof of Askey's conjectured $q$-analogue of Selberg's integral
and a conjecture of Morris,
\emph{SIAM J.\ Math.\ Anal.}\ \textbf{19}(4) (1988), 969--986.

\bibitem{klw-2015}
G.~K\'arolyi, A.~Lascoux, and S.~O.~Warnaar,
Constant term identities and Poincar\'e polynomials,
\emph{Trans.\ Amer.\ Math.\ Soc.}\ \textbf{367}(10) (2015),
6809--6836.
arXiv:1209.0855.

\bibitem{macdonald-1982}
I.~G.~Macdonald,
Some conjectures for root systems,
\emph{SIAM J.\ Math.\ Anal.}\ \textbf{13}(6) (1982), 988--1007.

\bibitem{mcnamara-2019}
P.~McNamara,
Positivity of cylindric skew Schur functions,
\emph{J.\ Combin.\ Theory Ser.~A}\ \textbf{165} (2019), 53--78.

\bibitem{morris-1982}
W.~G.~Morris,
\emph{Constant Term Identities for Finite and Affine Root Systems},
Ph.D.\ thesis, University of Wisconsin--Madison, 1982.

\bibitem{postnikov-2005}
A.~Postnikov,
Affine approach to quantum Schubert calculus,
\emph{Duke Math.\ J.}\ \textbf{128}(3) (2005), 473--509.
arXiv:math/0205165.

\bibitem{rosas-zabrocki}
M.~H.~Rosas and M.~Zabrocki,
Symmetric functions and the Kronecker product,
unpublished manuscript.

\bibitem{selberg-1944}
A.~Selberg,
Remarks on a multiple integral,
\emph{Norsk Mat.\ Tidsskr.}\ \textbf{26} (1944), 71--78.

\bibitem{warnaar-qselberg}
S.~O.~Warnaar,
$q$-Selberg integrals and Macdonald polynomials,
\emph{Ramanujan J.}\ \textbf{10}(2) (2005), 237--268.

\bibitem{warnaar-zudilin-2011}
S.~O.~Warnaar and W.~Zudilin,
A $q$-rious positivity,
\emph{Aequat.\ Math.}\ \textbf{81}(1--2) (2011), 177--183.

\bibitem{zeilberger-1982}
D.~Zeilberger,
A combinatorial proof of Dyson's conjecture,
\emph{Discrete Math.}\ \textbf{41}(3) (1982), 317--321.

\bibitem{zeilberger-bressoud-1985}
D.~Zeilberger and D.~M.~Bressoud,
A proof of Andrews' $q$-Dyson conjecture,
\emph{Discrete Math.}\ \textbf{54}(2) (1985), 201--224.

\bibitem{zeilberger-koh}
D.~Zeilberger,
Kathy O'Hara's constructive proof of the unimodality of the
Gaussian polynomials,
\emph{Amer.\ Math.\ Monthly}\ \textbf{96}(7) (1989), 590--602.

\end{thebibliography}

\end{document}
